\documentclass[autodetect-engine,ja=standard,b4j,landscape,twoside]{bxjsarticle}
%% [↓ ここから(オプション以外)触らない部分]===================
\usepackage{silence}
\WarningFilter{latex}{You have requested package}
\usepackage[sep,score]{sty/exam}
%% [↑ ここまで(オプション以外)触らない部分]===================

%% [↓ ここから設定] =============================
%% 模範解答を表示する場合は下のコメントを外す
\模範解答モードtrue

%% 再試の場合は下のコメントを外す
%\再試モードtrue
%% 追試の場合は下のコメントを外す
%\追試モードtrue

% 右上の欄の表示／非表示
%\欠課時数欄表示フラグtrue
\試験得点欄表示フラグtrue
\成績欄表示フラグtrue

\試験学年{4}
\試験クラス{I}
\試験タイトル{○年 サンプル科目 前期期末試験}
\試験日{2022年2月1日}
%\試験日{\qquad 年\qquad 月\qquad 日}

\setlength{\parindent}{0pt} % 字下げしない
%\graphicspath{{./fig/}} % 画像置き場
\renewcommand{\baselinestretch}{1.1} % 行間
\setlength{\arrayrulewidth}{1pt} % 表の線の太さ
\renewcommand{\labelenumi}{(\arabic{enumi}) } % 順序あり箇条書きの見出し
%% [↑ ここまで設定] =============================


%% よく使う自作コマンドをここに登録しましょう．
\newcommand{\ska}[1]{\left(#1\right)}
\newcommand{\cka}[1]{\left\{#1\right\}}
\newcommand{\dka}[1]{\left[#1\right]}

\begin{document}
\makeexamheader[6mm]

\vspace*{-5zh}
\underline{\bf 全ての問題で部分点はなしとする。}


\begin{mondai}[穴埋め問題のサンプル (20点 : 10点 $\times$ 2)]
以下の空欄に最も適切な語句を語群から選び、空欄に書け．
同じ語句を複数回選んでも良い．

%% 記入しやすいように少し行間を広げて，字下げしてます．
\vspace{1zh}
\renewcommand{\baselinestretch}{1.3}\selectfont
\setlength{\parindent}{1zw} % 字下げ

穴埋め問題はこのように作成します．
例えば 2 を 3 倍すると \穴抜き{6} になります．
もしカッコの広さが足りなかったり広すぎる場合，
\穴抜き[12zw]{このようにすることで} サイズを調整できます．

もし語群が必要であれば，以下のようにします．
中央寄せが不要なら centering 命令を消してください．
語群の中の語句はそのまま記載しても良いですが，
mboxに入れることで語句の途中での改行を防ぐことができます．

\begin{語群}
\centering
\mbox{りんご}，
\mbox{みかん}，
\mbox{ぶどう}，
\mbox{いちご}，
\mbox{メロン}，
\mbox{すいか}，
\mbox{なし}，
\mbox{もも}，
\mbox{いちじく}，
\mbox{グレープフルーツ}，
\mbox{パイナップル}，
\mbox{カレー}
\end{語群}
\end{mondai}


\begin{mondai}[計算問題のサンプル (30点 : 10点 $\times$ 3)]
計算問題などはこのように作成します．
\begin{enumerate}
  \setlength{\itemsep}{2zh}
    \item $1 + 2 \times 3$
      \回答欄{7}

    \item $\dfrac{\sigma_x}{2} \dfrac{\sigma_y}{3}$
      \回答欄{$\dfrac{\sigma_z}{6}$}

    \item たろう君はこの前の試験の点数をじろう君に見せました．
    じろう君はA君に「たろう君の2倍の点数を取った」と自慢し，
    B君には「たろう君があと1点取っていれば，たろう君には負けていた」と話しています．
    じろう君が正直者である場合，たろう君の点数は何点でしょう．
    ただし，点数は0以上100以下の整数とする．\\[1zh]
      \回答欄{0}
\end{enumerate}
\end{mondai}


\newpage



\begin{mondai}[表埋め問題のサンプル (20点 : 10点 $\times$ 2)]
表を埋める問題はこのように作成します．
\vspace{1zh}

\begin{tabular}{|l|l|}
\hline
\rowcolor{pgray}\bf 説明 & \bf 物の名前\\\hline
電話もできるアラーム機能付きゲーム機 & \答え{スマートフォン} \\\hline
とても素晴らしい文書作成ソフト & \答え{\LaTeX} \\\hline
\end{tabular}
\end{mondai}




\begin{mondai}[画像挿入サンプル (0点)]
  \includegraphics[width=0.5\linewidth]{fig/fig.png}
\end{mondai}



\裏面アナウンス
\newpage


\begin{mondai}[記述問題のサンプル (30点)]
記述問題はこのようにします．
\vspace{1zh}

\回答領域{
ここに長文の回答が入ります．
ここに長文の回答が入ります．
ここに長文の回答が入ります．
ここに長文の回答が入ります．
ここに長文の回答が入ります．
ここに長文の回答が入ります．
ここに長文の回答が入ります．
ここに長文の回答が入ります．
}

\end{mondai}

\end{document}
